\chapter{Primary-Source Catalog}
\label{v4-arith:app:primary-sources}

This appendix collects the primary-literature references that
anchor the arithmetic branch of Volume~IV. Each entry records
author(s), title, year, arxiv~ID (where applicable), and the role
in the branch (which chapter uses it for what).

\section{Kapranov and Kapranov--Smirnov}
\label{v4-arith:app:A:kapranov}

\begin{description}
\item[\textbf{Kapranov--Smirnov 1995}, \emph{Cohomology determinants
and reciprocity laws: number field case}.] Unpublished preprint,
$\sim$15 pp., ca.~1995. PDF mirrored at
\texttt{http://cage.ugent.be/\char`~kthas/Fun/library/KapranovSmirnov.pdf}.
Frames $\mathrm{Spec}\,\mathbb{Z}$ as an ``arithmetic curve'' via
the analogy between number fields and function fields over finite
fields; motivates the absolute point / $\mathbb{F}_{1}$ programme
and prime-labelled cohomology determinants feeding reciprocity laws.
Used throughout this branch as the foundational informal reference
(Chapters~\ref{v4-arith:drinfeld-kapranov}, \ref{v4-arith:arakelov-bar}).
\item[\textbf{Kapranov ca.~1995--96}, \emph{Analogies between number
fields and 3-manifolds}.] Unpublished preprint. Foundational for
arithmetic topology; the MKR (Mazur--Kapranov--Reznikov) dictionary
matches closed oriented 3-manifolds to number fields and knots to
primes. Extended in Reznikov arXiv:math/0107210 (2001) and
Morishita's book \emph{Knots and Primes} (2012). Used in
Chapter~\ref{v4-arith:gaiotto} for the Wilson-line $=$ ramified
Hecke character dictionary in Kim's arithmetic Chern--Simons.
\item[\textbf{Kapranov--Vasserot 2004}, \emph{Vertex algebras and
the formal loop space}, Publ.~IHES~100:209--269;
arXiv:math/0107143.] Introduces the ind-scheme $\mathcal{L}X$ of
formal loops into a finite-type scheme $X$; factorization structure
on $\mathcal{L}X$ recovers vertex-algebra structure on chiral de
Rham. Four-paper series; Part~IV = arXiv:math/0612371. No
$\mathrm{Spec}\,\mathbb{Z}$ version exists; this branch is the
arithmetic analogue.
\item[\textbf{Kapranov 2000}, \emph{The elliptic curve in the
S-duality theory and Eisenstein series for Kac--Moody groups},
arXiv:math/0001005.] S-duality on an elliptic curve; not used here
directly.
\end{description}

\section{Borger and $\Lambda$-Ring Theory}
\label{v4-arith:app:A:borger}

\begin{description}
\item[\textbf{Borger 2008}, \emph{The basic geometry of Witt
vectors, I: The affine case}, Algebra \& Number Theory 5:231--285;
arXiv:0801.1691.] Concrete description of \'etale algebras over
Witt vectors via commuting Frobenius lifts. Defines $\Lambda$-ring,
$\Lambda$-module distinctions used in
Chapter~\ref{v4-arith:gelfand-borger}.
\item[\textbf{Borger 2010}, \emph{The basic geometry of Witt
vectors, II: Spaces}, arXiv:1006.0092.] Extends to algebraic spaces;
generalises Buium's arithmetic jet functor. Used in
Chapter~\ref{v4-arith:gelfand-borger}~\S7 on Witt vectors.
\item[\textbf{Borger 2009}, \emph{$\Lambda$-rings and the field
with one element}, arXiv:0906.3146.] Central reference. Wilkerson
criterion (Prop.~1.2): $\Lambda$-ring = commuting $\psi_{p}$ with
$\psi_{p}(x)\equiv x^{p}\pmod{pR}$. Slogan (\S0 and \S2):
$\Lambda$-structure is descent from $R$ to
$R\otimes_{\mathbb{F}_{1}}\mathbb{Z}$. Used throughout
Chapter~\ref{v4-arith:gelfand-borger}.
\item[\textbf{Borger--de Smit 2008}, \emph{Galois theory and
integral models of $\Lambda$-rings}, arXiv:0801.2352.] Relative
case: $\Lambda$-structure over a Noetherian base.
\end{description}

\section{Bost--Connes, Connes--Marcolli, Connes--Consani}
\label{v4-arith:app:A:connes-etc}

\begin{description}
\item[\textbf{Bost--Connes 1995}, \emph{Hecke algebras, type III
factors and phase transitions with spontaneous symmetry breaking in
number theory}, Selecta Math.~1:411--457.] Central reference for
the arithmetic quantum statistical system. KMS partition function
$\zeta(\beta)$; phase transition at $\beta=1$; spontaneous
symmetry breaking by $\mathrm{Gal}(\mathbb{Q}^{\mathrm{ab}}/\mathbb{Q})$.
Explicitly motivated by Julia's primon gas.
\item[\textbf{Connes--Marcolli 2004}, \emph{From Physics to Number
Theory via Noncommutative Geometry, Part~I~\& II},
arXiv:math/0404128; arXiv:math/0411114.] $HH^{\bullet}(C^{*}_{\mathbb{Q}})$
computed \S1.8; used in Chapter~\ref{v4-arith:connes-soibelman}
for the Theorem~H$^{\mathrm{arith}}$ comparison.
\item[\textbf{Connes--Consani 2014}, \emph{The Arithmetic Site},
arXiv:1405.4527.] Arithmetic site $=$ Grothendieck topos
$\widehat{\mathbb{N}^{\times}}$ with tropical semiring
$\mathbb{Z}_{\max}$ structure sheaf. Used in
Chapter~\ref{v4-arith:gelfand-borger}~\S8 for the correct
$\mathbb{F}_{1}$-descent framework.
\item[\textbf{Connes--Consani 2015}, \emph{Geometry of the
arithmetic site}, arXiv:1502.05580.] Extends (1405.4527); Frobenius
correspondences as subvarieties.
\item[\textbf{Connes--Consani 2012}, \emph{Universal thickening of
the field of real numbers}, arXiv:1202.4377.]
\end{description}

\section{Kontsevich, Soibelman, Davison: $E_{n}$ Formality and CoHA}
\label{v4-arith:app:A:kontsevich-etc}

\begin{description}
\item[\textbf{Kontsevich 2003}, \emph{Deformation quantization of
Poisson manifolds}, Lett.~Math.~Phys., arXiv:q-alg/9709040.]
Foundational for operadic formality. Used in
Chapter~\ref{v4-arith:kontsevich-soibelman}~\S1.
\item[\textbf{Kontsevich 1999}, \emph{Operads and Motives in
Deformation Quantization}, arXiv:math/9904055.] $E_{n}$-formality
over $\mathbb{Q}$ for all $n$.
\item[\textbf{Tamarkin 1998}, \emph{Another proof of M.~Kontsevich
formality theorem}, arXiv:math/9803025.] Drinfeld-associator-based
proof for $E_{2}$.
\item[\textbf{Kontsevich--Soibelman 2010}, \emph{Cohomological Hall
algebra, exponential Hodge structures, and motivic DT invariants},
arXiv:1006.2706.] CoHA framework; used for
Chapter~\ref{v4-arith:connes-soibelman}~\S6 Soibelman-motivic
reading of Kesten--McKay.
\item[\textbf{Fresse 2017}, \emph{Homotopy of Operads and
Grothendieck--Teichm\"uller Groups}, Vol.~II, arXiv:1603.02939.]
Integral $E_{1}$-formality theorem used in
Chapter~\ref{v4-arith:kontsevich-soibelman}~\S1.
\item[\textbf{Kaledin 2014}, \emph{Universal Witt vectors and the
``Japanese cocycle''}, arXiv:1409.1859.] Integral-$E_{2}$
obstruction.
\item[\textbf{Salvatore--Wahl 2002}, arXiv:math/0202092.]
Dyer--Lashof operations on $H_{\bullet}(E_{2};\mathbb{F}_{p})$.
\item[\textbf{Davison 2016, 2017}, arXiv:1601.02479; arXiv:1705.01775.]
BPS Lie algebra PBW theorem; used for arithmetic CoHA conjecture in
Chapter~\ref{v4-arith:kontsevich-soibelman}~\S5.
\end{description}

\section{Beilinson--Drinfeld, Francis--Gaitsgory, Positselski}
\label{v4-arith:app:A:beilinson-etc}

\begin{description}
\item[\textbf{Beilinson--Drinfeld 2004}, \emph{Chiral Algebras},
AMS Colloq.~Publ.~51.] Ch.~3 (chiral bar), Ch.~4 (Hochschild).
\item[\textbf{Francis--Gaitsgory 2011}, \emph{Chiral Koszul
duality}, Selecta Math., arXiv:1103.5803.]
\item[\textbf{Positselski 2011}, \emph{Two kinds of derived
categories, Koszul duality, and comodule-contramodule
correspondence}, Mem.~AMS, arXiv:0905.2621.] Used in
Chapter~\ref{v4-arith:arakelov-bar}~\S7 for arithmetic Positselski.
\item[\textbf{Parshin, Kato (1979--1980).}] \emph{A generalization of local class field
theory by using K-groups.} Parshin reciprocity on the Gersten
complex. Used in
Chapter~\ref{v4-arith:arakelov-bar}~\S2 for Arnold-relation
replacement.
\end{description}

\section{Costello, Costello--Gwilliam, Costello--Gaiotto}
\label{v4-arith:app:A:costello}

\begin{description}
\item[\textbf{Costello--Gwilliam 2017, 2021}, \emph{Factorization
Algebras in Quantum Field Theory}, CUP Vol.~I \& II.] Used
throughout Chapter~\ref{v4-arith:costello-gwilliam}.
\item[\textbf{Costello 2013}, \emph{Supersymmetric Gauge Theory and
the Yangian}, arXiv:1303.2632.] Holomorphic-topological factorization.
\item[\textbf{Costello 2020}, \emph{Holography and Koszul duality},
arXiv:2006.10839.]
\item[\textbf{Costello--Gaiotto 2019}, \emph{Vertex operator
algebras and 3d $\mathcal{N}=4$ gauge theories}, arXiv:1804.06460.]
Used in Chapter~\ref{v4-arith:gaiotto}~\S3.
\item[\textbf{Costello--Dimofte--Gaiotto 2020}, \emph{Boundary
chiral algebras and holomorphic twists}, arXiv:2005.00083.]
\item[\textbf{Costello--Witten--Yamazaki 2017}, \emph{Gauge theory
and integrability I}, arXiv:1709.09993.] 4D Chern--Simons.
\item[\textbf{Costello--Yamazaki 2019}, \emph{Gauge theory and
integrability III}, arXiv:1908.02289.]
\item[\textbf{Francis 2012}, \emph{Factorization homology of
topological manifolds}, arXiv:1206.5522.] Used in
Chapter~\ref{v4-arith:costello-gwilliam}~\S4 for locally-constant
$E_{n}$-reduction.
\end{description}

\section{Geometric Langlands, Derived Satake, Hecke Categories}
\label{v4-arith:app:A:langlands}

\begin{description}
\item[\textbf{Mirkovi\'c--Vilonen 2007}, \emph{Geometric Langlands
duality and representations of algebraic groups over commutative
rings}, Ann.~Math., arXiv:math/0401222.] Geometric Satake.
\item[\textbf{Bezrukavnikov 2016}, \emph{On two geometric
realizations of an affine Hecke algebra}, Publ.~IHES,
arXiv:1209.0403.] Iwahori Hecke category; central for
Chapter~\ref{v4-arith:hochschild-trace}.
\item[\textbf{Ben-Zvi--Nadler 2009}, \emph{The character theory of
a complex group}, arXiv:0904.1247.] Hochschild trace $=$
$\mathcal{O}(\mathrm{LocSys})$; gluing principle for the synthesis
(Chapter~\ref{v4-arith:synthesis}).
\item[\textbf{Arinkin--Gaitsgory 2015}, \emph{Singular support of
coherent sheaves and the geometric Langlands conjecture}, Selecta
Math., arXiv:1201.6343.] Nilpotent singular support AG-conjecture.
\item[\textbf{Fargues--Scholze 2021}, \emph{Geometrization of the
local Langlands correspondence}, arXiv:2102.13459.] Local
arithmetic Langlands via Fargues--Fontaine curve.
\item[\textbf{Caraiani--Scholze 2017}, \emph{On the generic part of
the cohomology of compact unitary Shimura varieties},
arXiv:1710.10543.]
\item[\textbf{Emerton--Gee 2020}, \emph{A geometric perspective on
the Breuil--M\'ezard conjecture for $GL_{2}$}, arXiv:2012.12403.]
Emerton--Gee stack of Galois deformations; central for
Chapter~\ref{v4-arith:hochschild-trace}~\S4.
\item[\textbf{Zhu 2020}, arXiv:2008.02998.] Coherent Springer theory.
\item[\textbf{Zhu 2016}, arXiv:1603.05593; Bhatt--Scholze 2016,
arXiv:1602.08693.] Witt-vector affine Grassmannian.
\end{description}

\section{Kim's Arithmetic Chern--Simons}
\label{v4-arith:app:A:kim}

\begin{description}
\item[\textbf{Kim 2015}, \emph{Arithmetic Chern--Simons theory~I},
arXiv:1510.05818.] Foundational definition.
\item[\textbf{Chung--Kim--Kim--Pappas--Park--Yoo 2017}, \emph{Abelian
arithmetic Chern--Simons theory and arithmetic linking numbers},
arXiv:1711.04801.] Used in Chapter~\ref{v4-arith:gaiotto}.
\item[\textbf{Morishita 2009}, \emph{Analogies between knots and
primes, 3-manifolds and number rings}, arXiv:0904.3399.] MKR dictionary.
\end{description}

\section{Hartnoll--Yang and Primon Gas}
\label{v4-arith:app:A:primon-gas}

\begin{description}
\item[\textbf{Hartnoll--Yang 2025}, \emph{The Conformal Primon Gas
at the End of Time}, arXiv:2502.02661.] Primary seed reference.
$L$-function as partition function of charged oscillators
labelled by primes. BKL--primon correspondence (thermodynamic).
\item[\textbf{Julia 1990}, \emph{Statistical theory of numbers}, in
\emph{Number Theory and Physics}, Springer Proc.~Phys.~47:276--293.]
Original primon gas: free Fock on primes, $Z(\beta)=\zeta(\beta)$.
\item[\textbf{Spector 1990}, \emph{Supersymmetry and the M\"obius
inversion function}, Commun.~Math.~Phys.~127:239--252.]
Supersymmetric primon gas; $1/\zeta(\beta)=\sum\mu(n)n^{-\beta}$.
\end{description}

\section{Macdonald, DAHA, Schiffmann--Vasserot}
\label{v4-arith:app:A:macdonald}

\begin{description}
\item[\textbf{Macdonald 1995}, \emph{Symmetric Functions and Hall
Polynomials}, 2nd ed., Oxford.] Ch.~I used throughout.
\item[\textbf{Cherednik 2005}, \emph{Double Affine Hecke Algebras},
CUP.] DAHA framework.
\item[\textbf{Etingof--Kirillov 1994}, \emph{Macdonald's polynomials
and representations of quantum groups}, arXiv:hep-th/9312064.]
\item[\textbf{Schiffmann--Vasserot 2011}, \emph{The elliptic Hall
algebra and the equivariant K-theory of the Hilbert scheme of
$\mathbb{A}^{2}$}, arXiv:0905.2555 and arXiv:1202.2756.] Central for
Chapter~\ref{v4-arith:etingof}.
\end{description}

\section{Nekrasov, AGT, Maulik--Okounkov}
\label{v4-arith:app:A:nekrasov}

\begin{description}
\item[\textbf{Nekrasov 2003}, \emph{Seiberg--Witten prepotential from
instanton counting}, ATMP, arXiv:hep-th/0206161.] $Z_{\mathrm{inst}}$
definition.
\item[\textbf{Nekrasov--Okounkov 2006}, \emph{Seiberg--Witten theory
and random partitions}, Progress Math., arXiv:hep-th/0306238.]
\item[\textbf{Alday--Gaiotto--Tachikawa 2010}, arXiv:0906.3219.]
AGT correspondence.
\item[\textbf{Maulik--Okounkov 2019}, \emph{Quantum groups and
quantum cohomology}, Ast\'erisque, arXiv:1211.1287.] $R$-matrix on
instanton moduli.
\item[\textbf{Nekrasov--Shatashvili 2009}, arXiv:0908.4052.] NS
limit.
\item[\textbf{MNOP 2003}, arXiv:math/0312059.] Plane partition
generating function.
\item[\textbf{Gaiotto 2009}, arXiv:0904.2715.] Class-$\mathcal{S}$.
\item[\textbf{Beem--Lemos--Liendo--Peelaers--Rastelli--van Rees 2015},
\emph{Infinite chiral symmetry in four dimensions},
Commun.~Math.~Phys., arXiv:1312.5344.] BLLPRvR VOA/Higgs branch.
\end{description}

\section{Segal, Frenkel--Ben-Zvi}
\label{v4-arith:app:A:segal-fbz}

\begin{description}
\item[\textbf{Segal 2002}, \emph{The definition of conformal field
theory}, in \emph{Topology, Geometry and QFT}, CUP, pp.~421--577.]
Segal axioms.
\item[\textbf{Frenkel--Ben-Zvi 2004}, \emph{Vertex Algebras and
Algebraic Curves}, 2nd ed., AMS.] Standard reference.
\item[\textbf{Huang 1997}, \emph{Two-Dimensional Conformal Geometry
and Vertex Operator Algebras}, Birkh\"auser.]
\item[\textbf{Bakalov--Kirillov 2001}, \emph{Lectures on Tensor
Categories and Modular Functor}, AMS.]
\item[\textbf{Deligne}, \emph{Le d\'eterminant de la cohomologie}.]
Determinant line bundle on arithmetic moduli.
\end{description}

\section{Tate's Thesis and Adelic Analysis}
\label{v4-arith:app:A:tate}

\begin{description}
\item[\textbf{Tate 1950}, \emph{Fourier analysis in number fields
and Hecke's zeta-functions}, Princeton PhD thesis; reprinted in
Cassels--Fr\"ohlich 1967 \emph{Algebraic Number Theory}, CUP.]
Adelic Poisson summation, functional equation as sewing axiom.
\item[\textbf{Weil 1967}, \emph{Basic Number Theory}, Springer.]
Restricted direct product, adelic framework.
\item[\textbf{Jacquet--Langlands 1970}, \emph{Automorphic Forms on
$GL(2)$}, SLN~114.]
\item[\textbf{Gelbart--Jacquet 2003}, arXiv:math/0312099.]
\item[\textbf{Gross 1998}, \emph{On the Satake isomorphism}, in
\emph{Galois Representations in Arithmetic Algebraic Geometry},
CUP, Prop.~3.6.]
\item[\textbf{Bump}, \emph{Automorphic Forms and Representations},
CUP 1998.] \S1.5, \S3.
\end{description}

\section{Drinfeld, Molev: Yangians and R-Matrices}
\label{v4-arith:app:A:yangians}

\begin{description}
\item[\textbf{Drinfeld 1985}, \emph{Hopf algebras and the quantum
Yang--Baxter equation}, Soviet Math.~Dokl.~32:254--258.]
Yang $R$-matrix and RTT.
\item[\textbf{Drinfeld 1988}, \emph{A new realization of Yangians
and quantized affine algebras}, Soviet Math.~Dokl.~36:212--216.]
\item[\textbf{Molev 2007}, \emph{Yangians and Classical Lie
Algebras}, AMS.] Standard reference.
\item[\textbf{Feigin--Odesskii 1997}, arXiv:alg-geom/9702015.]
Shuffle algebras.
\item[\textbf{Frenkel--Reshetikhin 1998}, arXiv:math/9810055.]
$q$-characters.
\item[\textbf{Kamnitzer--Pham--Weekes 2018}, arXiv:1811.08779.]
Bethe algebra of opers.
\end{description}

\section{Arakelov, Deligne, Arithmetic Height Pairings}
\label{v4-arith:app:A:arakelov}

\begin{description}
\item[\textbf{Arakelov 1974}, \emph{Theory of intersection on
arithmetic surfaces}, Izv.~Akad.~Nauk SSSR 38.]
\item[\textbf{Faltings 1984}, \emph{Calculus on arithmetic
surfaces}, Ann.~Math.~119.]
\item[\textbf{Gillet--Soul\'e 1992}, \emph{Arithmetic intersection
theory}, Publ.~IHES~72.] Used in Chapter~\ref{v4-arith:nekrasov} for
Arakelov Chern classes.
\item[\textbf{Durov 2007}, arXiv:0704.2030.] $\mathbb{F}_{1}$-projective
space.
\item[\textbf{Deninger 1994}, \emph{Motivic $L$-functions and
regularized determinants}, Proc.~Symp.~Pure Math.~55.1; earlier
arXiv:alg-geom/9307001.]
\item[\textbf{Szpiro 1985}, Ast\'erisque~127.] Arakelov
$K_{\overline{\mathrm{Spec}\,\mathbb{Z}}}$ has degree $\log(2\pi)$.
\item[\textbf{Bost--K\"unnemann 2008}, arXiv:0806.4586.] Arakelov
vector bundles.
\end{description}

\section{Kesten--McKay}
\label{v4-arith:app:A:kesten-mckay}

\begin{description}
\item[\textbf{Kesten 1959}, \emph{Symmetric random walks on groups},
TAMS~92.]
\item[\textbf{McKay 1981}, \emph{The expected eigenvalue
distribution of a large regular graph}, Linear Algebra Appl.~40.]
\item[\textbf{Lubotzky 1994}, \emph{Discrete Groups, Expanding
Graphs and Invariant Measures}, Birkh\"auser, Thm~7.4.5.] $GL_{2}(\mathbb{Q}_{p})$
interpretation.
\item[\textbf{Terras 2011}, \emph{Zeta Functions of Graphs}, CUP.]
Ihara zeta; tree-moment identification used in
Chapter~\ref{v4-arith:etingof}.
\end{description}

\section{Prismatic Cohomology}
\label{v4-arith:app:A:prismatic}

\begin{description}
\item[\textbf{Bhatt--Scholze 2019}, \emph{Prisms and prismatic
cohomology}, arXiv:1905.08229.] Cited in
Chapter~\ref{v4-arith:costello-gwilliam}~\S3 as the upgrade path for
Fontaine crystalline data that might provide a non-trivial
$p$-adic diagonal thickening.
\end{description}

\section{Additional Cross-Volume References}
\label{v4-arith:app:A:cross-volume}

\begin{description}
\item[Vol~I (this repository)] CLAUDE.md; chapters of
\texttt{/Users/raeez/chiral-bar-cobar/}; five-theorem structure;
landscape census.
\item[Vol~II \texttt{/Users/raeez/chiral-bar-cobar-vol2/}]
$\mathsf{SC}^{\mathrm{ch,top}}$, 3D HT QFT, 3D quantum gravity.
\item[Vol~III \texttt{/Users/raeez/calabi-yau-quantum-groups/}]
Two-stage factorisation $\Phi_{d}=\mathrm{Sp}^{\mathrm{ch}}\circ\Phi^{\mathrm{FA}}_{d}$;
Borcherds/BKM $\Phi_{N}=c_{N}(0)/2$.
\end{description}
